Los experimentos fueron ejecutados en un entorno Linux
sobre WSL (Windows Subsystem for Linux), utilizando un
procesador AMD Ryzen 3 7320U with Radeon Graphics (2.40 GHz),
Sistema operativo de 64 bits, procesador x64 y 8 GB de memoria
RAM. Los programas principales fueron implementados en C++17
y compilados con \texttt{g++} utilizando el nivel de
optimización predeterminado vía \texttt{Makefile}. 

\vspace{0.4cm}

La medición de tiempo se realizó con la librería estándar
\texttt{<chrono>} (\texttt{high\_resolution\_clock}),
mientras que el consumo de memoria \textit{Peak RSS} se
registró mediante la función \texttt{getrusage()} de la
cabecera POSIX \texttt{<sys/resource.h>}. 

\vspace{0.4cm}

Cada medición corre en un proceso independiente para que 
el pico no se arrastre entre algoritmos. El conjunto de 
pruebas son 288 instancias generadas con un script de python 
con tamaños $n \in \{4, 8, 16, 64, 256, 1024\}$, 
$m \in \{n, 2n\}$, cuatro terrenos (libre, murallas, minas, mixta), 
dos dominios (D10, D1000) y tres semillas. La fuerza bruta 
solo se ejecutó con $N+M \le 32$. Cada punto de las Figuras 3 y 4 
promedia sobre dominios y semillas; en las Figuras 1 y 2 cada 
punto promedia solo las 3 semillas de un dominio. 

\vspace{0.4cm}

Las visualizaciones graficas fueron generadas mediante
scripts de Python usando las libreria \texttt{pandas} y
\texttt{seaborn}.